% ===== Opening: title, hook, the book and the series, today's plan (~5 min) =====

\webinartitle{4}{Գլուխ 4. Կարծում ես՝ խելացի՞ ես}%
  {\chip{numbers}{IQ} \chip{jeans}{պրոֆեսորներ} \chip{loudspeaker}{վստահություն} \chip{puzzle}{մարզում} \chip{giraffe}{հասակ}}%
  {smug}{Խելքը չի հաջողվում չափել}{{graduation}{chess}{exploding}}

% Hook: five everyday "intelligence glitches"; the audience types the ones they know
\begin{frame}{Ծանո՞թ է}
  \begin{columns}[c]
    \begin{column}{0.66\textwidth}
      \begin{popup}{ուղեղ.exe}
        {\blackfont Խելքը չի արձագանքում}\par\smallskip
        \begin{enumerate}\small
          \item<2-> Համացանցի թեստն ասաց՝ IQ-դ 160 է
          \item<3-> Գերազանցիկ ընկերդ ձու խաշել չգիտի
          \item<4-> Ոչինչ չկարդացածն ամենավստահն է
          \item<5-> Լավ պատրաստվել էիր, բայց համոզված ես, որ կտրվել ես
          \item<6-> «Ուղեղի մարզման» խաղում ռեկորդ ունես, քննությանը՝ ոչ մի օգուտ
        \end{enumerate}
      \end{popup}
    \end{column}
    \begin{column}{0.3\textwidth}
      \centering
      \mascot[2.4cm]{confused}\par\medskip
      \onslide<7->{%
      \begin{chatbox}
        Գրիր քեզ ծանոթ համարները (օր.՝ \textbf{2, 4})
      \end{chatbox}}
    \end{column}
  \end{columns}
\end{frame}

% The book and where this chapter sits in the series (for newcomers)
\begin{frame}{Գիրքը և շարքը}
  \begin{columns}[T]
    \begin{column}{0.5\textwidth}
      \begin{card}
        {\blackfont\Large The Idiot Brain} {\small (2016)}\par\smallskip
        \textbf{Դին Բըրնեթ}՝ նեյրոգիտնական Քարդիֆի համալսարանից ու ստենդափ կատակերգու \emoji{grin}\par\smallskip
        Գաղափարը՝ ուղեղն ապշեցուցիչ է, բայց \textbf{հեռու է կատարյալ լինելուց}։
      \end{card}
      \vspace{1mm}
      \begin{card}[colback=mint!35!cardwhite]
        Նախորդները բաց ես թողե՞լ։ Ոչինչ. ամեն գլուխ առանձին պատմություն է։
      \end{card}
    \end{column}
    \begin{column}{0.44\textwidth}
      {\blackfont\small Մեկ գլուխ ամեն վեբինարում}\par\medskip
      {\color{plum!55!cream}\small \#1. Ներածություն, գլուխ 1\par\smallskip
      \#2. Գլուխ 2. Հիշողություն\par\smallskip
      \#3. Գլուխ 3. Վախ}\par\smallskip
      \chip[butter]{brain}{\#4. Գլուխ 4. Խելք (այսօր)}\par\smallskip
      {\color{plum!55!cream}\small \#5. Գլուխ 5. Ընկալում\par\smallskip
      \#6. Գլուխ 6. Անհատականություն\par\smallskip
      \#7. Գլուխ 7. Ուրիշ մարդիկ\par\smallskip
      \#8. Գլուխ 8. Երբ ուղեղը «փչանում» է}
    \end{column}
  \end{columns}
\end{frame}

% Plan for today
\begin{frame}{Այսօր}
  \begin{columns}[T]
    \begin{column}{0.47\textwidth}
      \begin{card}[colback=peach!40!cream]
        {\blackfont Մինչև ընդմիջումը}\par\medskip
        \chip{numbers}{IQ}\par\smallskip
        Ինչ է չափում IQ թեստը\par\medskip
        \chip{jeans}{Պրոֆեսորի տաբատը}\par\smallskip
        Խելքի տեսակները, ու ինչու են խելացիները հիմարություններ անում
      \end{card}
    \end{column}
    \begin{column}{0.47\textwidth}
      \begin{card}[colback=lavender!45!cream]
        {\blackfont Ընդմիջումից հետո}\par\medskip
        \chip{loudspeaker}{Վստահություն}\ Ինչու են քիչ գիտեցողներն ավելի վստահ\par\smallskip
        \chip{puzzle}{Ուղեղի մարզում}\ Կարելի՞ է խելացի դառնալ\par\smallskip
        \chip{giraffe}{Հասակ ու գեներ}\ Ժառանգականություն ու մշակույթ
      \end{card}
    \end{column}
  \end{columns}
  \vspace{3mm}
  \centering
  {\small \emoji{coffee}\ Մեջտեղում՝ 5 րոպե ընդմիջում \qquad \emoji{speech}\ Հարցերը գրիր չաթում ցանկացած պահի}
\end{frame}
